\documentclass[11pt,a4paper]{article}
\usepackage[margin=2.5cm]{geometry}
\usepackage[T1]{fontenc}
\usepackage{fancyhdr}   % headers and footers
\usepackage{booktabs}   % professional table rules
\usepackage{array}      % extra column types (m{}, b{})

% ---------------------------------------------------------------
% MODULE II - Styling Pages
% Syllabus units 5-7 (Text 1, Tutorial VI and VII).
% Compiles with a single run of pdflatex.
% ---------------------------------------------------------------

\pagestyle{fancy}
\fancyhf{}                                  % clear default header/footer
\fancyhead[L]{MAT5FS112 -- Module II}
\fancyhead[R]{\thepage}
\renewcommand{\headrulewidth}{0.4pt}

\title{Module II: Styling Pages}
\author{Study Notes -- MAT5FS112}
\date{}

\begin{document}
\maketitle

% ===============================================================
\section{Fonts and text styling}
% ===============================================================

Font commands change the look of a piece of text:
\begin{center}
\begin{tabular}{ll}
\toprule
Command & Result \\
\midrule
\verb|\textbf{bold}| & \textbf{bold} \\
\verb|\textit{italic}| & \textit{italic} \\
\verb|\emph{emphasis}| & \emph{emphasis} (italic inside italic text becomes upright) \\
\verb|\texttt{typewriter}| & \texttt{typewriter} \\
\verb|\textsf{sans serif}| & \textsf{sans serif} \\
\verb|\textsc{small caps}| & \textsc{small caps} \\
\verb|\underline{underlined}| & \underline{underlined} \\
\bottomrule
\end{tabular}
\end{center}

Use \verb|\emph| for meaning and \verb|\textit| for pure styling.
Commands can be combined: \verb|\textbf{\textit{both}}|.

\subsection{Font size}

Sizes are set with declarations. A declaration applies until the end of
the enclosing group (a pair of braces or an environment):
\begin{verbatim}
{\small This is small text.}
{\Large This is large text.}
\end{verbatim}

Sizes, from small to large: \verb|\tiny|, \verb|\scriptsize|,
\verb|\footnotesize|, \verb|\small|, \verb|\normalsize|, \verb|\large|,
\verb|\Large|, \verb|\LARGE|, \verb|\huge|, \verb|\Huge|.

{\small This is small text.} \quad {\Large This is large text.}

% ===============================================================
\section{Spacing and page layout}
% ===============================================================

\subsection{Paragraph spacing}

\begin{verbatim}
\setlength{\parindent}{0pt}     % no indent at the start of paragraphs
\setlength{\parskip}{0.6em}     % vertical gap between paragraphs
\noindent                       % suppress indent for one paragraph
\end{verbatim}

\subsection{Manual spacing}

\verb|\vspace{1cm}| adds vertical space and \verb|\hspace{2cm}| adds
horizontal space. \verb|\\| forces a line break and \verb|\newpage|
forces a page break.

\subsection{Margins}

The \texttt{geometry} package sets margins in one line, as in this file:
\begin{verbatim}
\usepackage[margin=2.5cm]{geometry}
\end{verbatim}

\subsection{Headers, footers and page numbers}

The \texttt{fancyhdr} package gives control over the running header
and footer, with \verb|L|, \verb|C| and \verb|R| positions. This
document uses:
\begin{verbatim}
\pagestyle{fancy}
\fancyhf{}                          % clear everything first
\fancyhead[L]{Left text}
\fancyhead[R]{\thepage}
\renewcommand{\headrulewidth}{0.4pt}
\end{verbatim}

Page numbering style: \verb|\pagenumbering{roman}| (i, ii, iii\ldots)
or \verb|\pagenumbering{arabic}|. \verb|\thispagestyle{empty}| removes
the header from one page.

% ===============================================================
\section{Displayed text}
% ===============================================================

Environments display blocks of text with a specific alignment or style.

\subsection{Alignment}

\begin{center}
Centred text.
\end{center}

\begin{flushleft}
Left-aligned (ragged right) text.
\end{flushleft}

\begin{flushright}
Right-aligned text.
\end{flushright}

\subsection{Quotations and verse}

\begin{quote}
A quotation: a paragraph set with extra margin on both sides.
\end{quote}

\begin{quotation}
The \texttt{quotation} environment is similar, but it indents the
first line of each paragraph inside it.
\end{quotation}

\begin{verse}
Verse keeps each line break you type\\
and centres nothing else on the page.
\end{verse}

\subsection{Side-by-side blocks with minipage}

\noindent
\begin{minipage}[t]{0.45\textwidth}
  Left column. A \texttt{minipage} is a box of fixed width that
  holds text. Its contents are wrapped inside the box.
\end{minipage}
\hfill
\begin{minipage}[t]{0.45\textwidth}
  Right column. Two minipages placed one after the other with
  \verb|\hfill| between them sit side by side.
\end{minipage}

\subsection{Verbatim text}

\begin{verbatim}
This text is printed exactly as typed:  $x^2$ and \textbf{stay raw}.
\end{verbatim}

% ===============================================================
\section{Rows and columns: tabular}
% ===============================================================

The \texttt{tabular} environment sets rows and columns. The argument
gives one letter per column: \texttt{l} (left), \texttt{c} (centre),
\texttt{r} (right), and \texttt{p\{width\}} (wrapped paragraph of
given width). Cells are separated by \verb|&| and rows end with
\verb|\\|.

\begin{verbatim}
\begin{tabular}{lcr}
Item   & Mark & Grade \\
\hline
Algebra & 18  & A     \\
Real    & 14  & B     \\
\end{tabular}
\end{verbatim}

\begin{tabular}{lcr}
Item    & Mark & Grade \\
\hline
Algebra & 18   & A     \\
Real    & 14   & B     \\
\end{tabular}

\subsection{Rules: \texttt{\textbackslash hline} and \texttt{\textbackslash cline}}

\verb|\hline| draws a full-width horizontal line. \verb|\cline{2-3}|
draws a line only under columns 2 to 3. Vertical lines come from
\texttt{|} in the column specification.

\begin{tabular}{|l|c|c|}
\hline
Name & Physics & Maths \\
\cline{2-3}
Asha & 80      & 75    \\
\hline
\end{tabular}

\subsection{Spanning columns: \texttt{\textbackslash multicolumn}}

\begin{tabular}{|l|c|c|}
\hline
\multicolumn{3}{|c|}{Semester V marks} \\
\hline
Name & Physics & Maths \\
\hline
Asha & 80      & 75    \\
\hline
\end{tabular}

\subsection{Professional tables with \texttt{booktabs}}

The \texttt{booktabs} package uses \verb|\toprule|, \verb|\midrule|,
\verb|\bottomrule| and avoids vertical lines, which gives cleaner
results. \verb|\cmidrule{2-3}| draws a partial rule.

\begin{tabular}{l r r}
\toprule
Module & Units & Min.\ marks \\
\midrule
I   & 4 & 10 \\
II  & 3 & 10 \\
III & 4 & 10 \\
\bottomrule
\end{tabular}

\subsection{Fixed-width columns with \texttt{p\{\}}}

Use \texttt{p\{width\}} when a cell has a paragraph of text. The
\texttt{array} package adds \texttt{m\{width\}} (vertically centred)
and \texttt{b\{width\}} (bottom aligned).

\begin{tabular}{p{3.5cm} p{8cm}}
\toprule
Command & Purpose \\
\midrule
\verb|\hline| & Draws a full horizontal rule across the table. \\
\verb|\multicolumn| & Merges several cells in a row into one. \\
\bottomrule
\end{tabular}

% ===============================================================
\section{Tables as floating objects}
% ===============================================================

A table placed in a \texttt{table} environment floats: LaTeX moves
it to a suitable place and numbers it. The caption and the label let
the text refer to it.

\begin{table}[htbp]
  \centering
  \caption{Marks awarded in the LaTeX internal examination}
  \label{tab:marks}
  \begin{tabular}{lcc}
    \toprule
    Student & Assignment & Viva \\
    \midrule
    Student A & 9 & 8 \\
    Student B & 7 & 9 \\
    \bottomrule
  \end{tabular}
\end{table}

As Table~\ref{tab:marks} shows, the marks are recorded in two parts.
The placement specifier \texttt{[htbp]} means: here, top, bottom,
or on its own page. Note the \verb|~| before \verb|\ref|: it is a
non-breaking space that keeps ``Table'' and the number together.

\end{document}
